% Notes-only candidate: a finite normalized-pair detector, not an outer-action theorem.
\subsection{A finite detector retaining both torsion frames}

\begin{statement}{C.47}{Open}{Candidate: the exact derived depth of the retained parity character}
Let $\mathcal H=\ker\lambda\subset\widehat{\mathrm{GT}}_0$.
Use its canonical normalized pairs and retain their conjugator $f$,
not only the induced outer action. For the reductions of these pairs
modulo $F^{(d)}$, the least $d$ through which $\Delta|_{\mathcal H}$
factors is
\[
                              d=4.                    \tag{C.47.1}
\]
In particular $f\in F^{(4)}$ implies $\Delta(1,f)=0$, while
C.46 supplies an actual full pair with $f\in F^{(3)}$ and
$\Delta(1,f)=1$.

There is an explicit finite characteristic quotient $C$ of
$F/F^{(4)}$ such that $\Delta|_{\mathcal H}$ factors through
the finite group of actual normalized pair residues on $C$.
The quotient and an exact program evaluating the character from
its retained residue are constructed below. The program uses a
base group of order $576$, coset modules of dimensions $288$ over
$\mathbb F_2$ and $192$ over $\mathbb F_3$, and finite linear
systems; it does not enumerate the full affine group.

The assertion concerns the retained normalized conjugator.
It does not identify the kernel of the outer action on
$F/F^{(4)}$, the smallest possible finite detector, the full
set of actual residues, or an explicit detector at all multipliers.
\end{statement}

\begin{proof}[Candidate proof, pending complete independent audits]
Use the temporary marking of C.44--C.46 on the actual pure subgroup
\[
 \Gamma=\widehat{C_2*C_3}=\langle s,r\mid s^2=r^3=1\rangle,
 \quad \tau=sr,\quad F=\ker(\Gamma\to S_3),\quad
 x=\tau^{-2},\quad y=sxs.
\]
For a multiplier-one normalized gentle pair, its actual extension
satisfies
\[
          \Phi(\tau)=\tau,\qquad \Phi(s)=sf,
          \qquad \Phi(r)=sf\tau.                     \tag{C.47.2}
\]
Its full pure torsion frames are defined by
$\Phi(r)=ArA^{-1}$ and $\Phi(s)=\beta s\beta^{-1}$,
with $A,\beta\in F$. These are the actual unique frames used in
C.20. The original marking is obtained by
$x_{\rm old}=y^{-1}$, $y_{\rm old}=x^{-1}$ and the corresponding
inner normalization; the normalized conjugator becomes $f^{-1}$
in the old letters. This change preserves every derived subgroup.
The frame and intertwiner change by the same inner adjustment,
so the CM coordinates and their difference character are unchanged.
One may therefore evaluate in the displayed temporary marking.

\paragraph{The finite base and its pure subgroup.}
Put
\[
 Q=\operatorname{SL}_2(\mathbb Z/12)/\{\pm I\},\quad
 S=\begin{pmatrix}0&1\\-1&0\end{pmatrix},\quad
 U_0=\begin{pmatrix}1&1\\0&1\end{pmatrix},\quad
 q(s)=[S],\quad q(r)=[SU_0^{-1}].                    \tag{C.47.3}
\]
Then $q(\tau)=[U_0^{-1}]$, and the signed pure generators are
$U_0^2$ and $SU_0^2S^{-1}$. Write $Q_0=q(F)$.
The exact finite groups have
\[
 |Q|=576,\qquad |Q_0|=96,\qquad
 |Q_0'|=8,\quad |Q_0''|=2,\quad Q_0'''=1.            \tag{C.47.4}
\]
Here $Q_0=\ker(Q\to\operatorname{SL}_2(\mathbb F_2))$.
The standalone certificate generates all the base matrices modulo
the single global sign and computes all these derived images.
In particular $Q_0$ meets $\langle q(s)\rangle$ and
$\langle q(r)\rangle$ trivially: their respective images in
$S_3$ have the same orders $2$ and $3$.

\paragraph{A finite centralizer lemma.}
Let a finite group $Q$ contain an element $h$ of prime order $p$.
Set $H=\langle h\rangle$, let $V=\mathbb F_p[Q/H]$ with its
left permutation action, and let $\delta_H$ denote the basis vector
of the identity coset. In $V\rtimes Q$ the element
\[
                       \widetilde h=(\delta_H,h)
\]
has exact order $p$: $h$ fixes $\delta_H$, so its norm is
$p\delta_H=0$. If $(a,b)$ centralizes $\widetilde h$, then
$b$ centralizes $h$ and the translation equation is
\[
                    (1-h)a=\delta_H-b\delta_H.       \tag{C.47.5}
\]
The coordinate of the left side at the $h$-fixed coset $H$ is zero.
The coordinate of the right side is $1$ unless $bH=H$.
Consequently $b\in H$. Conversely the powers of
$\widetilde h$ project onto $H$. Thus
\[
       \operatorname{im}\bigl(C_{V\rtimes Q}(\widetilde h)\to Q\bigr)
                      =\langle h\rangle.            \tag{C.47.6}
\]
Adding any other $Q$-module does not change this obstruction in
the displayed coordinate. No cohomology vanishing or splitting
assertion is being substituted for the literal equation.

\paragraph{Simultaneously retaining the two frames.}
For the $Q$ of (C.47.3), form
\[
 V_s=\mathbb F_2[Q/\langle q(s)\rangle],\qquad
 V_r=\mathbb F_3[Q/\langle q(r)\rangle],\qquad
 V=V_s\oplus V_r,\qquad G=V\rtimes Q.
\]
Their dimensions are $288$ and $192$. Choose actual affine lifts
\[
 \widetilde s=(\delta_{\langle s\rangle},0,q(s)),\qquad
 \widetilde r=(0,\delta_{\langle r\rangle},q(r)).     \tag{C.47.7}
\]
They have orders $2$ and $3$, respectively. The universal property
of the profinite free product gives a continuous homomorphism
$\pi:\Gamma\to G$ taking $s,r$ to these elements. Put
$\widetilde\tau=\widetilde s\widetilde r$.
Its pure restriction lands in the explicitly defined finite group
\[
                   L=V\rtimes Q_0,
             \qquad |L|=2^{293}3^{193}.              \tag{C.47.8}
\]
Surjectivity of $\pi$ onto $G$ or of $\pi(F)$ onto $L$ is not
needed. Because $V$ is abelian and $Q_0'''=1$, one has
\[
                              L^{(4)}=1.             \tag{C.47.9}
\]

For any actual full gentle pair, the single residue $\pi(f)$
determines
\[
 s_f=\widetilde s\,\pi(f),\qquad
 r_f=s_f\widetilde\tau.                              \tag{C.47.10}
\]
There are actual conjugators $\pi(\beta),\pi(A)\in L$ taking
$\widetilde s,\widetilde r$ to $s_f,r_f$. If two elements of
$L$ conjugate either torsion lift to the same target, their quotient
centralizes that lift. By (C.47.6) and
$Q_0\cap\langle q(s)\rangle=Q_0\cap\langle q(r)\rangle=1$,
their projections to $Q$ agree. Therefore
\[
                       q(\beta),\ q(A)
       \quad\text{are uniquely determined by }\pi(f).           \tag{C.47.11}
\]
This argument does not require $\Phi$ to descend to an automorphism
of $G$: the identities are the images of the actual full frame
identities under $\pi$.

\paragraph{An exact finite evaluation algorithm.}
Represent an element of $L$ by its two translation vectors and
its base matrix in $Q_0$. For $h=\widetilde s$ or
$\widetilde r$, and its target in (C.47.10), enumerate the
$96$ possible base conjugators $b\in Q_0$. Once the base
conjugacy equation holds, the translation equation is
\[
           (1-k)a=w-bv,                              \tag{C.47.12}
\]
where $(v,h_0)$ is the initial lift and $(w,k)$ the target.
This is a linear equation in each of the two permutation modules.
For each cycle of $k$, it has a solution precisely when the sum
of the right-hand coordinates on that cycle is zero; choosing one
coordinate zero and successive partial sums supplies a solution.
Existence is guaranteed for actual full residues, and (C.47.11)
proves uniqueness of the resulting base. The algorithm may solve
in all of $L$, since any additional finite solution has the same
base as the actual one.

C.12 then recovers the normalized unipotent intertwiner
$V=\left(\begin{smallmatrix}1&t\\0&1\end{smallmatrix}\right)$
modulo twelve by testing $t=0,\ldots,11$ against both base images
of $s,r$. There is exactly one such $t$ for an actual pair;
the direct matrix equations also show that no nonidentity such
unipotent centralizes both base generators. Put
\[
                      T=q(A)^{-1}V,\qquad
                      W=q(\beta)^{-1}V.             \tag{C.47.13}
\]
For each of $T$ and $W$ separately, choose one global sign,
used consistently at both moduli three and four, to obtain its
genuine lift equal to $I$ modulo four, as in C.20 and
R24 E.65. This choice exists and is unique for the actual CM
coordinates. At three, $T$ lies in the order-six cubic centralizer
and $W$ in the order-four quartic centralizer. If their signed
modulo-three matrices are $T_3,W_3$, then
\[
 T_3^3=(-I)^{\epsilon_3},\qquad
 W_3^2=(-I)^{\epsilon_4},\qquad
             \Delta=\epsilon_3+\epsilon_4\pmod2.     \tag{C.47.14}
\]
These are precisely the cubic sign and the quartic class modulo
$\{\pm I\}$ in R24 E.65. Separate primewise sign choices are
not made.

\paragraph{A characteristic finite pair quotient.}
Define
\[
 K=\bigcap_{\theta\in\operatorname{Hom}_{\rm cts}(F,L)}\ker\theta,
                       \qquad C=F/K.                \tag{C.47.15}
\]
There are exactly $|L|^2$ maps, since $F$ is free profinite on
two generators. Thus $C$ is finite, is fully characteristic,
and has order at most $|L|^{|L|^2}$. By (C.47.9), every map
kills $F^{(4)}$, so $F^{(4)}\subset K$. The fixed map
$\pi|_F$ factors through this $C$.

For normalized multiplier-one pairs,
$\alpha_f(x)=x$ and $\alpha_f(y)=f^{-1}yf$, and the exact
product is $f_{gh}=f_g\alpha_g(f_h)$. A fully characteristic $K$ is
preserved by every $\alpha_g$, and its induced action on $C$
is determined by the retained $f_g\bmod K$. Hence reduction
is a homomorphism onto a finite group of actual pair residues.
The algorithm (C.47.10)--(C.47.14) depends only on the projection
of that residue to the one specified $L$-factor. It proves that
$\Delta|_{\mathcal H}$ factors through this finite group.
It does not assert that every finite-compatible pair is an
actual residue or enumerate the full residue image.

\paragraph{The sharp retained derived depth.}
If $f\in F^{(4)}$, then $\pi(f)=1$. Equation (C.47.10)
fixes the two affine torsion generators. The unique pure bases
in (C.47.11) are therefore both one, and the intertwiner is one.
Thus $T=W=1$ and $\Delta=0$. Equivalently, the character kills
the kernel of normalized pair reduction modulo $F^{(4)}$.
C.46 supplies an actual full pair in the kernel of reduction
modulo $F^{(3)}$ with value one. It lies in the kernels at all
smaller derived levels too. This proves (C.47.1).
No claim about the outer-action kernel at the fourth derived
level follows without an additional centralizer argument forcing
its normalized conjugator into $F^{(4)}$.

The executed standalone program
\path{code/pbc_finite_parity_detector.py}
constructs the complete finite base and coset permutations,
checks the centralizer-coordinate obstruction for every base
centralizer element, verifies the affine torsion relations and
pure derived orders, and implements (C.47.12)--(C.47.14).
It checks the identity and a common pure inner-frame change.
A separate finite affine test exercises parity one; it is explicitly
labelled a finite frame test and not a full lifting result.
The actual nonzero full witness is C.46. The initial program and
outputs before that additional test are also retained.
\end{proof}
